\section{Signaux inter-threads en espace utilisateur}

Dans la version finale de la bibliothèque, nous avons ajouté un mécanisme de
signaux inter-threads entièrement en espace utilisateur, sans recourir aux
signaux système POSIX. Ce choix s'imposait car le signal \texttt{SIGVTALRM}
est déjà utilisé par le mode préemptif, et utiliser des signaux système pour
la communication entre threads aurait introduit des interférences difficiles à
contrôler.

\subsection{Représentation des signaux}

Les signaux sont représentés par un masque de bits \texttt{uint32\_t}, où le
bit $i$ indique que le signal numéro $i$ est présent. Chaque thread dispose de
trois nouveaux champs dans sa structure \texttt{thread\_s} :

\begin{itemize}
    \item \texttt{pending\_signals} : masque des signaux reçus mais pas encore
    traités.
    \item \texttt{sig\_handlers[32]} : tableau de handlers, un par numéro de
    signal (\texttt{NULL} si aucun handler n'est enregistré).
    \item \texttt{sigwait\_mask} : masque des signaux qu'un thread attend
    explicitement via \texttt{thread\_sigwait}. Ces signaux ne sont pas livrés
    aux handlers tant que le thread est en attente.
\end{itemize}

\subsection{Fonctions implémentées}

Trois fonctions ont été ajoutées, disponibles uniquement en mode
\texttt{USE\_CONTEXT} :

\begin{itemize}
    \item \texttt{thread\_kill(thread\_t t, int sig)} : pose le bit du signal
    \texttt{sig} dans \texttt{pending\_signals} du thread cible. Si ce thread
    est bloqué dans un \texttt{thread\_sigwait} attendant précisément ce
    signal, il est immédiatement remis dans la \texttt{readyQueue}.

    \item \texttt{thread\_signal(int sig, handler)} : enregistre un handler
    pour le signal \texttt{sig} du thread courant et retourne l'ancien handler,
    à la manière de l'appel système \texttt{signal()}.

    \item \texttt{thread\_sigwait(uint32\_t set, int *signo)} : bloque le
    thread courant jusqu'à la réception d'un signal appartenant au masque
    \texttt{set}. Si un signal correspondant est déjà en attente au moment de
    l'appel, la fonction retourne immédiatement sans bloquer.
\end{itemize}

\subsection{Livraison coopérative}

La bibliothèque étant coopérative, les signaux ne sont pas livrés de façon
asynchrone. La fonction interne \texttt{deliver\_pending\_signals()} est
appelée aux points de rendez-vous naturels : au retour du \texttt{swapcontext}
dans \texttt{thread\_yield}, et au démarrage de chaque thread dans
\texttt{thread\_wrapper}. Elle parcourt \texttt{pending\_signals \&
\textasciitilde{}sigwait\_mask} et invoque les handlers enregistrés, en
excluant les signaux réservés à un \texttt{thread\_sigwait} en cours.

\subsection{Test : \texttt{91-signals.c}}

Nous avons ajouté le test \texttt{91-signals.c} qui couvre trois scénarios :

\begin{enumerate}
    \item \textbf{Sigwait bloquant :} un thread appelle \texttt{thread\_sigwait}
    et se bloque. Un second thread lui envoie le signal via
    \texttt{thread\_kill}. Le premier thread doit se réveiller et recevoir le
    bon numéro de signal.

    \item \textbf{Handler de signal :} un thread enregistre un handler pour un
    signal, se l'envoie à lui-même, puis appelle \texttt{thread\_yield}. Le
    handler doit être appelé exactement une fois lors du yield.

    \item \textbf{Signal déjà en attente :} le signal est envoyé avant que le
    receveur n'appelle \texttt{thread\_sigwait}. La fonction doit alors
    retourner immédiatement, le signal étant déjà présent dans
    \texttt{pending\_signals}.
\end{enumerate}

Le test se compile via \texttt{make signals}, qui utilise
\texttt{-DUSE\_CONTEXT} et lie contre \texttt{libthread-context.a}.
